\section{زنجیره بلوکی و \lr{IOTA Tangle}}

\subsection{مرور کلی بر زنجیره بلوکی}

در فصل ۴ مشاهده شد که وابستگی به مرجع قابل اعتماد\LTRfootnote{\lr{Trusted Authority (TA)}} به‌عنوان نقطه شکست منفرد\LTRfootnote{\lr{Single Point of Failure (SPoF)}}، امنیت کل شبکه را در معرض خطر قرار می‌دهد. از سوی دیگر، مدیریت اعتماد و شهرت که در فصل ۷ بررسی شد، برای کار خود به بستری نیاز دارد که امتیازات اعتماد و سوابق رفتار گره‌ها را به‌صورت ایمن، تغییرناپذیر و قابل اعتماد ذخیره کند. زنجیره بلوکی\LTRfootnote{\lr{blockchain}} به‌عنوان راهکاری برای پاسخ به این نیازها معرفی می‌شود: شبکه‌ای همتا به همتا که اجرای ایمن تراکنش‌ها را بدون نیاز به شخص ثالث مورد اعتماد ممکن می‌سازد \cite{razaBlockchainBasedReputationTrust2024}. در این ساختار، هر بلوک شامل سرآیند، گواهی و هش بلوک قبلی است و زنجیره‌ای پیوندی رو به عقب از بلوک‌ها ایجاد می‌کند؛ سرآیند هر بلوک شامل شناسهٔ بلوک، شناسهٔ تولیدکننده، امضای بلوک، هش بلوک و مهر زمانی است و گواهی بلوک نیز شناسهٔ گواهی، شناسهٔ گره، کلید عمومی گره، هش گواهی و مهر زمانی را در خود دارد \cite{razaBlockchainBasedReputationTrust2024}. آنچه باعث تغییرناپذیری این زنجیره می‌شود، گنجاندن مداوم هش بلوک قبلی در هر بلوک جدید است؛ بلوکی که در دفتر کل ثبت شود، به‌صورت مقاوم در برابر دست‌کاری و در سرتاسر شبکه تکثیر می‌شود \cite{razaBlockchainBasedReputationTrust2024}. در مدل معرفی‌شده در مقالهٔ \lr{SBTMS}، سرآیند بلوک شامل شاخص، هش قبلی، تعداد تراکنش‌ها، مهر زمانی، نانس و درخت مرکل است و زنجیره بلوکی دفتر کلی توزیع‌شده و مبتنی بر اجماع است که همهٔ تراکنش‌های موفق در فهرستی از بلوک‌ها ذخیره می‌شود \cite{ghovanlooyghajarSBTMSScalableBlockchain2021}. زنجیره بلوکی بر پایهٔ فناوری دفتر کل توزیع‌شده\LTRfootnote{\lr{Distributed Ledger Technology (DLT)}} بنا شده و شبکهٔ همتا به همتا، دفتر کل مشترک، رمزنگاری و بستر محاسباتی را در هم ترکیب می‌کند \cite{razaBlockchainBasedReputationTrust2024}.

معماری زنجیره بلوکی در مرور رضا و همکاران از شش لایهٔ اصلی تشکیل شده است (این طرح شش‌لایه در آن مرور به دو منبع اولیه ارجاع شده است) \cite{razaBlockchainBasedReputationTrust2024}. لایهٔ داده ساختار بلوک‌ها و زنجیره، توابع هش، مهر زمانی، درخت مرکل و امضای دیجیتال را در بر می‌گیرد؛ لایهٔ شبکه شبکهٔ همتا به همتا، پروتکل‌های انتقال و سازوکار تأیید را شامل می‌شود؛ لایهٔ توافق مکانیزم‌هایی مانند اثبات کار، اثبات سهم، اثبات سهم نمایندگی و تحمل خطای بیزانس عملی را جای می‌دهد؛ لایهٔ انگیزش سازوکارهای صدور و توزیع ارز را مدیریت می‌کند؛ لایهٔ قرارداد شامل کدهای اسکریپت، قراردادهای هوشمند و الگوریتم‌ها و سازوکارهای مرتبط است؛ و لایهٔ کاربرد در نهایت دانش مالی قابل برنامه‌نویسی، ارز قابل برنامه‌نویسی و جامعهٔ قابل برنامه‌نویسی را پوشش می‌دهد \cite{razaBlockchainBasedReputationTrust2024}.

\begin{figure}[H]
\centering
\includegraphics[width=0.8\textwidth]{blockchainstruct1.png}
\caption{ساختار عمومی زنجیره بلوکی: زنجیره‌ای از بلوک‌ها با اتصال هش هر بلوک به بلوک قبلی}
\label{fig:blockchain_struct}
\end{figure}

شکل \ref{fig:blockchain_struct} ساختار عمومی زنجیره بلوکی و نحوهٔ اتصال هر بلوک به هش بلوک قبلی را نشان می‌دهد. از حیث دسترسی نیز زنجیره‌های بلوکی به سه نوع مجوزدار (خصوصی)، بدون مجوز (عمومی) و کنسرسیومی (ترکیبی) تقسیم می‌شوند؛ در نوع مجوزدار فقط اعضای یک گروهٔ بسته حق نوشتن دارند، در نوع بدون مجوز هر تراکنش برای همهٔ گره‌ها قابل مشاهده است و هر گره می‌تواند در اجماع شرکت کند، و نوع کنسرسیومی نیمه‌خصوصی است و به گروهی از سازمان‌های تأییدشده داده می‌شود \cite{razaBlockchainBasedReputationTrust2024}.

\subsubsection{مکانیزم‌های توافق}

پس از مرور ساختار زنجیره بلوکی، سازوکار اجماع را که مهم‌ترین مؤلفهٔ لایهٔ توافق است بررسی می‌کنیم. الگوریتم اجماع راه‌حلی است برای مسئلهٔ ژنرال‌های بیزانسی و کار آن تعیین این است که یک بلوک جدید برای درج در دفتر کل مشروع است یا نه؛ الگوریتم کارا باید در برابر حمله مقاوم، تحمل‌خطا، کارسا و در دسترس همهٔ شرکت‌کنندگان شبکه باشد \cite{razaBlockchainBasedReputationTrust2024}. این الگوریتم‌ها را می‌توان به دو دستهٔ مبتنی بر اثبات و مبتنی بر رأی‌گیری تقسیم کرد \cite{razaBlockchainBasedReputationTrust2024}: در دستهٔ نخست گره‌ای که تأیید کافی انجام دهد، امتیاز افزودن بلوک جدید را می‌یابد و پاداش می‌شود، و در دستهٔ دوم اجماع بر پایهٔ تحمل خطای بیزانس و با اثباتات ریاضی محکم بنا می‌شود \cite{razaBlockchainBasedReputationTrust2024}.

در اثبات کار\LTRfootnote{\lr{Proof of Work (PoW)}} که نخستین پروتکل اجماع زنجیره بلوکی است و برای مقابله با اسپم توسط درک و نائور پیشنهاد شد، اعتبارسنج باید پازل محاسباتی حل کند؛ این الگوریتم نیاز به توان محاسباتی بالا دارد، مقیاس‌پذیری ضعیفی دارد و تأخیر طولانی برای تأیید تراکنش ایجاد می‌کند \cite{razaBlockchainBasedReputationTrust2024}. در اثبات سهم\LTRfootnote{\lr{Proof of Stake (PoS)}} اعتبارسنج به‌صورت شبه‌تصادفی بر اساس سهام اقتصادی همتا در شبکه انتخاب می‌شود و قطعیت بلوک در آن سریع‌تر از اثبات کار است، هرچند اعتبارسنج‌های با سهام بزرگ می‌توانند به تمرکززدایی، کاهش امنیت در مقایسه با اثبات کار و فقدان قطعیت اجماع منجر شوند \cite{razaBlockchainBasedReputationTrust2024}. اثبات سهم نمایندگی\LTRfootnote{\lr{Delegated Proof of Stake (DPoS)}} انرژی کمتری مصرف می‌کند و تراکنش‌ها را سریع‌تر از اثبات کار و اثبات سهم انجام می‌دهد و از سازوکار «یک رأی به ازای هر سهم» بهره می‌گیرد، اما غیرمتمرکزسازی را محدود می‌کند \cite{razaBlockchainBasedReputationTrust2024}. تحمل خطای بیزانس عملی\LTRfootnote{\lr{Practical Byzantine Fault Tolerance (PBFT)}} در سال ۱۹۹۹ برای کاهش پیچیدگی الگوریتم اصلی تحمل خطای بیزانس از نمایی به چندجمله‌ای معرفی شد و پرکاربردترین الگوریتم اجماع است، هرچند برای برخی کاربردها مانند مدیریت زنجیرهٔ تأمین و معاملات انرژی مناسب نیست \cite{razaBlockchainBasedReputationTrust2024}.

\begin{table}[H]
\centering
\caption{مکانیزم‌های توافق در زنجیره بلوکی}
\label{tab:consensus}
\begin{tabular}{p{2.5cm} p{4cm} p{2.5cm} p{2.5cm}}
\toprule
\textbf{الگوریتم} & \textbf{ویژگی کلیدی} & \textbf{مصرف انرژی} & \textbf{مقیاس‌پذیری} \\
\midrule
\lr{PoW} & پازل محاسباتی؛ نخستین پروتکل اجماع & بالا & پایین \\
\lr{PoS} & سهام اقتصادی؛ انتخاب شبه‌تصادفی اعتبارسنج & پایین & متوسط \\
\lr{DPoS} & یک رأی به ازای هر سهم & بسیار پایین & بالا \\
\lr{PBFT} & کاهش پیچیدگی بیزانسی از نمایی به چندجمله‌ای (۱۹۹۹) & پایین تا متوسط & بالا (در شبکه‌های بزرگ ناکارآمد) \\
\bottomrule
\end{tabular}
\end{table}

جدول \ref{tab:consensus} ویژگی‌های کلیدی این چهار الگوریتم را خلاصه می‌کند. نکتهٔ قابل توجه تناقض درونی منبع دربارهٔ \lr{PBFT} است: جدول مقایسهٔ الگوریتم‌های مبتنی بر رأی‌گیری مقیاس‌پذیری آن را «بالا» گزارش می‌کند، اما در همان ردیف محدودیت «عدم کارآمدی در شبکه‌های بزرگ» را ذکر می‌کند؛ تحمل مهاجم این الگوریتم نیز حداکثر برابر یک‌سوم گره‌های معیوب (حدود ۳۳ درصد) است و برای شبکه‌های بزرگ ناکارآمد است و در برابر حملات ممانعت از سروی و سیبل آسیب‌پذیر است \cite{razaBlockchainBasedReputationTrust2024}. انتخاب الگوریتم اجماع نیز به الزامات کاربرد خاص از جمله امنیت، مقیاس‌پذیری و کارایی انرژی بستگی دارد \cite{razaBlockchainBasedReputationTrust2024}.

\subsection{چالش‌های زنجیره بلوکی در شبکه خودرویی}

با اینکه زنجیره بلوکی ویژگی‌های امنیتی ذاتی مانند یکسانی، مقاومت در برابر دست‌کاری، شبه‌نام‌پذیری و مقاومت در برابر چندین حمله دارد \cite{razaBlockchainBasedReputationTrust2024}، برای به‌کارگیری در شبکه‌های خودرویی با چالش‌هایی روبه‌روست. به‌دلیل فرایند ماینینگ، بیشتر فناوری‌های زنجیره بلوکی از قصوری مانند توان عملیاتی پایین و مقیاس‌پذیری ضعیف رنج می‌برند \cite{ghovanlooyghajarSBTMSScalableBlockchain2021}، درحالی‌که تعداد و سرعت بالای تراکنش‌های تولیدشده در شبکه‌های خودرویی مدام در حال افزایش است \cite{ghovanlooyghajarSBTMSScalableBlockchain2021}. پروتکل‌های سنتی مانند اثبات کار محاسباتی سنگین و کند هستند و تأخیر بالایی ایجاد می‌کنند که برای شبکه‌های پویا مانند شبکه‌های خودرویی نامناسب است و علاوه بر این، اتریوم و هایپرلج توانایی محدودی در تعداد تراکنش در ثانیه دارند \cite{razaBlockchainBasedReputationTrust2024}. به‌عنوان مثال، تعداد تراکنش در ثانیه حدود ۷ برای بیت‌کوین و حدود ۱۵ برای اتریوم گزارش شده است (در مقالهٔ مورد نظر منبع اولیهٔ این ارقام ذکر نشده است) \cite{IOTATANGLE20}. نویسندگان دیگر این مسئله را «سه‌گانهٔ مقیاس‌پذیری» می‌نامند: زنجیره‌های بلوکی می‌توانند هم‌زمان فقط دو از سه ویژگی غیرمتمرکز، امن و مقیاس‌پذیر را داشته باشند \cite{silvanoIotaTangleCryptocurrency2020}.

مصرف انرژی بالای اثبات کار چالش دیگری است؛ این الگوریتم مصرف انرژی بالا و هزینهٔ محاسباتی سنگین دارد \cite{razaBlockchainBasedReputationTrust2024} و الزامات توان و مصرف انرژی بالا و کار محاسباتی پرهزینهٔ اثبات کار از مشکلات طراحی بلاکچین برای سناریوهای ناهمگون دستگاه‌ها به شمار می‌رود \cite{silvanoIotaTangleCryptocurrency2020}. نگرانی‌های انرژی ناشی از فرایند ماینینگ نیز بسیاری از راه‌حل‌های بلاکچینی را برای کاربردهای \lr{IoT} نامناسب می‌کند \cite{sealeyIOTATangle202022}. به این چالش‌ها، کارمزد تراکنش نیز افزوده می‌شود؛ تراکنش‌های زنجیره بلوکی کارمزد دارند \cite{IOTATANGLE20, silvanoIotaTangleCryptocurrency2020} و وجود کارمزد، زمان پردازش بالا و فقدان مقیاس‌پذیری در سناریوهایی با دستگاه‌های ناهمگون جای نمی‌گیرد \cite{silvanoIotaTangleCryptocurrency2020}. افزون بر این، قیمت و کارمزد نوسانی ارزهای دیجیتال می‌تواند هزینه‌های غیرقابل پیش‌بینی ایجاد کند \cite{silvanoIotaTangleCryptocurrency2020}.

تأخیر بالا و حجم ذخیره‌سازی نیز از محدودیت‌های این فناوری هستند؛ تأخیر تعهد تراکنش و نیاز به ذخیره‌سازی حجم بزرگی از داده‌ها از چالش‌های پیش رو هستند \cite{silvanoIotaTangleCryptocurrency2020} و «ظرفیت ذخیره‌سازی» نیز در فهرست چالش‌های مرور رضا و همکاران عنوان شده است \cite{razaBlockchainBasedReputationTrust2024}. از حیث حریم خصوصی، زنجیره بلوکی محرمانگی را در میان اهداف اصلی خود ندارد و از دست رفتن کلید خصوصی می‌تواند به دوبارخرج‌کردن یا وارونگی تراکنش منجر شود \cite{razaBlockchainBasedReputationTrust2024}. از حیث امنیت، حملات به کیف پول و حملهٔ ۵۱ درصد شناخته‌شده‌ترین تهدیدهای امنیتی زنجیره بلوکی هستند؛ در حملهٔ ۵۱ درصد گروهی از ماینرهای بدخواه بیش از ۵۰ درصد نرخ هش شبکه را تصاحب می‌کنند و می‌توانند دوبارخرج‌کردن و بازآرایی دفتر کل انجام دهند \cite{razaBlockchainBasedReputationTrust2024}.

چالش‌های باقی‌مانده بیشتر سازمانی و سازگاری‌محور هستند. گذار از مدل متمرکز به غیرمتمرکز تلاش فنی زیادی می‌خواهد و یکی از چالش‌های فوری، ادغام سامانهٔ مدیریت اعتماد در سامانه‌های موجود است \cite{razaBlockchainBasedReputationTrust2024}؛ دفترهای کل توزیع‌شده می‌توانند برای تعامل‌پذیری میان ذی‌نفعان متعدد به کار روند، اما این کار پیچیدگی زمانی اضافی ایجاد می‌کند \cite{farsimadanReviewSecurityChallenges2025}. سایر چالش‌های عنوان‌شده شامل پیچیدگی طراحی، حکمرانی و مقررات، مسئلهٔ انشعاب و پیچیدگی محاسباتی است \cite{razaBlockchainBasedReputationTrust2024}. الگوریتم‌های سنتی اجماع نیاز به توان محاسباتی بالا، شکل‌گیری اجماع کند و توان عملیاتی داده پایین دارند \cite{razaBlockchainBasedReputationTrust2024} و حتی مدل‌هایی که خود را غیرمتمرکز می‌نامند، اغلب همچنان به شخص ثالث مورد اعتماد یا مرجع صدور گواهی نیاز دارند \cite{razaBlockchainBasedReputationTrust2024}. در بافت خودرویی، ورود داده نادرست عملیات شبکه را تحت تأثیر قرار می‌دهد و به انتشار خطا می‌انجامد \cite{razaBlockchainBasedReputationTrust2024} و دستیابی به اجماع میان صاحبان خودرو و نهادهای دولتی برای نظارت و هماهنگی دشوار است \cite{razaBlockchainBasedReputationTrust2024}.

\subsection{مزایا و کاربردهای زنجیره بلوکی در شبکه اقتضایی خودرویی}

پس از بررسی چالش‌ها، ویژگی‌هایی هستند که زنجیره بلوکی را برای شبکه‌های اقتضایی خودرویی جذاب می‌کند. غیرمتمرکزسازی با حذف نقطه شکست منفرد، امنیت و تاب‌آوری را افزایش می‌دهد \cite{razaBlockchainBasedReputationTrust2024}؛ در سامانه‌های متمرکز اگر سامانهٔ اصلی از کار بیفتد کل شبکه از کار می‌افتد، در حالی که در بلاکچین نیازی به شخص ثالث مورد اعتماد برای برقراری اعتماد نیست \cite{ghovanlooyghajarSBTMSScalableBlockchain2021}. تغییرناپذیری تراکنش‌ها تضمین می‌کند که داده‌ها قابل تغییر یا دست‌کاری نباشند و لایه‌ای اضافی از اعتماد فراهم کنند \cite{razaBlockchainBasedReputationTrust2024}، و ردیابی‌پذیری تراکنش‌ها از سوءاستفادهٔ مخربانه از داده‌ها جلوگیری می‌کند \cite{razaBlockchainBasedReputationTrust2024}. تراکنش‌ها با شفافیت کامل قابل ردیابی هستند چون زنجیره بلوکی تاریخچهٔ کامل تراکنش‌ها را نگه می‌دارد \cite{razaBlockchainBasedReputationTrust2024}. قراردادهای هوشمند می‌توانند بدون وابستگی به یک شخص ثالث مورد اعتماد، سیاست‌های کنترل دسترسی را اجرا کنند \cite{razaBlockchainBasedReputationTrust2024} و برای ارزیابی اعتماد خودروها نیز به‌کار رفته‌اند؛ مثلاً دو قرارداد \lr{FASC} و \lr{TESC} برای ارزیابی اعتماد خودروها پیشنهاد شده‌اند \cite{razaBlockchainBasedReputationTrust2024}. افزون بر این، ویژگی‌های غیرمتمرکز، توزیع‌شده، ذخیره‌سازی گسترده و عدم قابلیت دست‌کاری از مشخصه‌های بلاکچین در این حوزه به شمار می‌رود \cite{ghovanlooyghajarSBTMSScalableBlockchain2021}.

در حوزهٔ شبکه‌های خودرویی، صحت پیام‌های پخش‌شده با ذخیرهٔ اعتماد در زنجیره بلوکی بررسی شده است؛ در طرح یانگ و همکاران خودروها اعتبار پیام دریافتی را رتبه‌بندی کرده و به واحد کنارجاده‌ای\LTRfootnote{\lr{Road-Side Unit (RSU)}} می‌فرستند و آن واحد آفست اعتماد را با اجماع ترکیبی اثبات کار و اثبات سهم محاسبه می‌کند \cite{razaBlockchainBasedReputationTrust2024}، و در طرح ژانگ و همکاران شهرت در زنجیره بلوکی ذخیره شده و پیام‌های نادرست شناسایی می‌شوند \cite{razaBlockchainBasedReputationTrust2024}. مقادیر اعتماد در بلوک‌ها ذخیره می‌شوند \cite{razaBlockchainBasedReputationTrust2024} و در برخی طرح‌ها مقادیر اعتماد خودروها در بلاکچین نگهداری و میان واحدهای کنارجاده‌ای به اشتراک گذاشته می‌شوند \cite{farsimadanReviewSecurityChallenges2025}. مدیریت هویت و کلید نیز کاربرد دیگری است؛ طرح تئودولی و همکاران با شناسه‌های غیرمتمرکز و اعتبارنامه‌های تأییدپذیر استاندارد \lr{W3C} از نقطه شکست منفرد اجتناب می‌کند و نیاز به مرجع صدور گواهی متمرکز را حذف می‌کند \cite{razaBlockchainBasedReputationTrust2024}، کلیدها در برخی طرح‌ها روی دفتر کل ذخیره می‌شوند \cite{razaBlockchainBasedReputationTrust2024} و در طرح دیگری تابع فیزیکی غیرقابل تکثیر\LTRfootnote{\lr{Physical Unclonable Function (PUF)}} با زیرساخت کلید عمومی بلاکچینی ترکیب شده است \cite{razaBlockchainBasedReputationTrust2024}. از حیث حریم خصوصی، در طرح لیو و همکاران نام مستعار (نشانی عمومی) در زنجیره بلوکی ثبت می‌شود، مرجع قابل اعتماد هویت مخرب را ردیابی می‌کند و امضای گروهی مبتنی بر هویت به‌کار می‌رود \cite{razaBlockchainBasedReputationTrust2024}؛ اثباتات دانش‌صفر و اثباتات دانش‌صفر غیرتعاملی نیز برای ناشناس‌ماندن و حریم خصوصی شرطی در شبکه‌های خودرویی گزارش شده است \cite{razaBlockchainBasedReputationTrust2024}.

دو طرح مبتنی بر بلاکچین برای شبکه‌های اقتضایی خودرویی نیز در این زمینه طراحی شده‌اند. سامانهٔ احراز هویت حافظ حریم خصوصی با کمک بلاکچین\LTRfootnote{\lr{Blockchain-Assisted Privacy-Preserving Authentication System (BPAS)}} احراز هویت را به‌صورت خودکار انجام می‌دهد و هم‌زمان حریم خصوصی خودرو را حفظ می‌کند؛ این طرح جز در دو مرحلهٔ راه‌اندازی سیستم و ثبت‌نام خودرو، به مرکز ثبت برخطی نیاز ندارد، امکان ردیابی شرطی و ابطال پویای خودروهای دارای رفتار نادرست را فراهم می‌کند و از نظر نویسندگان کارا و مقیاس‌پذیر است \cite{fengBPASBlockchainAssistedPrivacyPreserving2020}. امنیت این طرح با تحلیل امنیتی عمیق و ارزیابی عملکردی جامع بر بستر هایپرلج فراب، یک زنجیرهٔ بلوکی مجوزدار، ارزیابی شده است \cite{fengBPASBlockchainAssistedPrivacyPreserving2020}. طرح \lr{BAIV}\LTRfootnote{\lr{BAIV}} نیز احراز هویت ناشنای کارآمد را با حفظ یکپارچگی پیام ترکیب می‌کند \cite{mariaBAIVEfficientBlockchainBased2022}؛ با ادغام زنجیره بلوکی در شبکهٔ اقتضایی خودرویی، تصدیق هویت کاربر خودرو بدون دخالت مرجع مورد اعتماد ممکن می‌شود و حریم خصوصی شرطی برای ابطال خودروهای مخرب در صورت اختلاف و جلوگیری از آسیب بیشتر به سامانه به‌کار می‌رود \cite{mariaBAIVEfficientBlockchainBased2022}. با گذر خودروی تصدیق‌شده از حوزهٔ یک واحد کنارجاده‌ای به حوزهٔ بعدی، نیاز به تصدیق مجدد از بین می‌رود: واحد کنارجاده‌ای رسید تصدیق هویت می‌سازد و آن را به واحدهای کنارجاده‌ای آینده می‌فرستد تا تصدیق مکرر اجتناب شود \cite{mariaBAIVEfficientBlockchainBased2022}. از نظر کارایی، هزینهٔ تصدیق هویت یک کاربر و یک واحد کنارجاده‌ای ۱۲٫۵۴ میلی‌ثانیه است (حدود ۸۰ کاربر در ثانیه و نزدیک به ۴۸۰۰ کاربر در دقیقه) و برای ۱۰۰ کاربر، ۱۲۹۱ میلی‌ثانیه در مقابل بیش از ۱۳۷۵ میلی‌ثانیه در چهار طرح پیشین گزارش شده است \cite{mariaBAIVEfficientBlockchainBased2022}. با این حال، مرجع مورد اعتماد همچنان برای ثبت‌نام آفلاین، نگاشت شناسه‌های ساختگی و ابطال لازم است، زنجیره بلوکی در ارزیابی عملکردی پیاده‌سازی نشده و تنها عملیات رمزنگاری زمان‌سنجی شده، و تحلیل امنیتی مقاله استدلالی است و اثبات صوری ارائه نمی‌دهد \cite{mariaBAIVEfficientBlockchainBased2022}.

\subsection{راهکارهای رفع محدودیت‌های زنجیره بلوکی}

برای رفع محدودیت‌های پیشین، راهکارهایی در منابع بررسی‌شده پیشنهاد شده است. شاردینگ\LTRfootnote{\lr{sharding}} به‌عنوان سازوکاری مطرح شده که با تقسیم حجم زیاد داده‌ها به تکه‌های کوچک‌تر، دفتر کل زنجیره بلوکی را کارآمدتر، مقیاس‌پذیرتر و پایدارتر می‌کند \cite{ghovanlooyghajarSBTMSScalableBlockchain2021}؛ الگوریتم شاردینگ وظایف ماینینگ را میان کمیته‌ها توزیع می‌کند و هر کمیته مجموعه‌ای از تراکنش‌ها را پردازش می‌کند \cite{ghovanlooyghajarSBTMSScalableBlockchain2021}. استفاده از الگوریتم‌های جایگزین اجماع نیز راهکاری است؛ به‌عنوان مثال، استفاده از اثبات سهم یا تحمل خطای بیزانس عملی می‌تواند تأخیر را کاهش دهد، اما ممکن است همچنان الزامات تأخیر پایین سخت‌گیرانهٔ کاربردهای بلادرنگ شبکه‌های خودرویی را برآورده نکند \cite{razaBlockchainBasedReputationTrust2024}. اثبات سهم همچنان مشکل مقیاس‌پذیری دارد که مانع استخراج تراکنش‌های خودرو به‌صورت نزدیک به بلادرنگ توسط واحدهای کنارجاده‌ای است \cite{ghovanlooyghajarSBTMSScalableBlockchain2021}، در حالی که ترکیب تحمل خطای بیزانس عملی و اثبات کار مقیاس‌پذیری را بهبود می‌بخشد و پیچیدگی محاسباتی را کاهش می‌دهد \cite{razaBlockchainBasedReputationTrust2024}. رویکردهای ترکیبی نیز شامل بهره‌گیری از نقاط قوت الگوریتم‌های شناخته‌شده یا ادغام فناوری‌های مختلف مانند فناوری ماینینگ مجازی است \cite{razaBlockchainBasedReputationTrust2024} و مدل نیمه‌متمرکز نیز بیشتر معایب هر دو مدل متمرکز و غیرمتمرکز را برطرف می‌کند \cite{razaBlockchainBasedReputationTrust2024}.

در ادامه مشاهده می‌شود که سامانهٔ \lr{SBTMS} از شاردینگ بهره می‌گیرد و نویسندگان آن تأکید می‌کنند برخلاف \lr{Tangle} که هماهنگ‌کننده‌اش نقطهٔ شکست منفرد است، الگوریتم شاردینگ کاملاً توزیع‌شده است \cite{ghovanlooyghajarSBTMSScalableBlockchain2021}. این مقایسه زمینه را برای معرفی \lr{IOTA Tangle} به‌عنوان جایگزینی با ساختار دفتر کل متفاوت فراهم می‌کند.

\subsubsection{سیستم مدیریت اعتماد مقیاس‌پذیر مبتنی بر زنجیره بلوکی \lr{(SBTMS)}}

سامانهٔ مدیریت اعتماد مقیاس‌پذیر مبتنی بر زنجیره بلوکی\LTRfootnote{\lr{Scalable Blockchain Trust Management System (SBTMS)}} یکی از طرح‌هایی است که با ترکیب شاردینگ و اجماع بیزانسی به محدودیت‌های مقیاس‌پذیری پاسخ می‌دهد \cite{ghovanlooyghajarSBTMSScalableBlockchain2021}. این سامانه یک بستر غیرمتمرکز ایجاد می‌کند و کار خود را در چهار گام انجام می‌دهد \cite{ghovanlooyghajarSBTMSScalableBlockchain2021}. در گام نخست، خودروی گیرنده امتیاز اطمینان رویداد را بر اساس فاصلهٔ فرستنده تا رویداد محاسبه می‌کند (پیام‌های خودروهای نزدیک‌تر به رویداد قابل اطمینان‌تر هستند) و قابلیت اطمینان پیام را با استنباط بیزی محاسبه کرده و نتیجه را به نزدیک‌ترین واحد کنارجاده‌ای می‌فرستد \cite{ghovanlooyghajarSBTMSScalableBlockchain2021}. در گام دوم، واحد کنارجاده‌ای «قابلیت اطمینان خالص» تجمیعی را با تابع سیگموید محاسبه می‌کند و تراکنش تنها زمانی تأیید می‌شود که این مقدار از آستانهٔ ۷۰ درصد بگذرد \cite{ghovanlooyghajarSBTMSScalableBlockchain2021}. در گام سوم، واحدهای کنارجاده‌ای برای جلوگیری از تبانی، به‌جای هویت دائمی یا زیرساخت کلید عمومی، هویت گذرای خودتولیدشده می‌سازند؛ آن‌ها با حل اثبات کار روی «تصادف دوره»\LTRfootnote{\lr{EpochRandomness}} شناسه‌ای گذرا تولید و به کمیته‌ها (شاردها) تقسیم می‌شوند \cite{ghovanlooyghajarSBTMSScalableBlockchain2021}. توجه داشته باشید که در این طرح اثبات کار حذف نشده بلکه فقط برای ساخت هویت گذرا و مقاومت در برابر تبانی و حملهٔ سیبل به‌کار رفته است \cite{ghovanlooyghajarSBTMSScalableBlockchain2021}. در گام چهارم، در هر کمیته پروتکل تحمل خطای بیزانس عملی با سه فاز پیش‌آماده‌سازی، آماده‌سازی و تعهد اجرا می‌شود و پذیرش پس از دریافت دو‌سوم پیام‌های آماده‌سازی صورت می‌گیرد؛ شارد اجماع‌شده با امضای حداقل نصف کمیته به‌علاوهٔ یک واحد کنارجاده‌ای به کمیتهٔ نهایی می‌رود و کمیتهٔ نهایی پس از محاسبهٔ خلاصهٔ رمزنگاری، دوباره \lr{PBFT} اجرا کرده و بلوک را منتشر می‌کند \cite{ghovanlooyghajarSBTMSScalableBlockchain2021}.

تحلیل امنیتی این طرح مدل تهدیدی با خودروها و واحدهای کنارجاده‌ای مخرب یا هک‌شده، پیام جعلی، گزارش جعلی قابلیت اطمینان و عوامل خودخواه در نظر می‌گیرد؛ استدلال می‌شود استنباط بیزی اثر پیام‌های جعلی را خنثی می‌کند و تصادف دوره ترکیب تصادفی کمیته را تضمین می‌کند، اما این تحلیل توصیفی است و اثبات صوری ارائه نمی‌دهد \cite{ghovanlooyghajarSBTMSScalableBlockchain2021}. در ارزیابی، شبیه‌سازی ناهمگام پایتون با ۲۰ واحد کنارجاده‌ای، ۱۰ تا ۳۰۰۰ خودرو، سه نوع رویداد (حادثه، ترافیک و تخریب جاده) و زمان شبیه‌سازی ۳۶۰۰ ثانیه انجام شده است \cite{ghovanlooyghajarSBTMSScalableBlockchain2021}. دقت و درستی نتایج با تعداد بلوک‌ها افزایش می‌یابد و الگوریتم شاردینگ عملکرد را تقریباً خطی با توان محاسباتی واحدهای کنارجاده‌ای مقیاس می‌کند \cite{ghovanlooyghajarSBTMSScalableBlockchain2021}. با افزایش تعداد خودروها، زمان تولید بلوک در این سامانه تقریباً خطی افزایش می‌یابد، در حالی که در اثبات کار با شیب بسیار بیشتری رشد می‌کند (بر اساس خوانش تقریبی شکل ۱۳ مقاله، حدود ۲ ثانیه در مقابل حدود ۱۴ ثانیه برای ۳۰۰۰ خودرو) و زمان محاسبهٔ خودرو نیز کمتر از اثبات کار است \cite{ghovanlooyghajarSBTMSScalableBlockchain2021}. یک ناسازگاری درونی مقاله نیز قابل ذکر است: جدول پارامترها بزرگی کمیته را ۵ گزارش می‌کند، در حالی که متن از کمیته‌هایی با اندازهٔ ۴ و سه واحد کنارجاده‌ای برای اجماع سخن می‌گوید \cite{ghovanlooyghajarSBTMSScalableBlockchain2021}.

\subsection{\lr{IOTA Tangle} به عنوان جایگزین}

درحالی‌که \lr{SBTMS} محدودیت‌های زنجیره بلوکی را با شاردینگ مهار می‌کند، \lr{IOTA Tangle} با تکیه بر ساختار دفتر کل، رویکردی متفاوت در برابر این چالش‌ها دارد. \lr{IOTA Tangle} یک فناوری دفتر کل توزیع‌شده است که برای شبکه‌ها و دستگاه‌های \lr{IoT} طراحی شده و بر پایهٔ ساختار گراف جهت‌دار بدون دور\LTRfootnote{\lr{Directed Acyclic Graph (DAG)}} بنا شده است \cite{sealeyIOTATangle202022, IOTATANGLE20, silvanoIotaTangleCryptocurrency2020}. در این ساختار بلوکی وجود ندارد؛ نویسندگان سیلوانو و مارچلینو زنجیره بلوکی را خود نوعی گراف جهت‌دار بدون دور محدود شده به یک مسیر (بلوک بعدی) می‌دانند و تنگل را گراف جهت‌دار بدون دور کمتر محدودی معرفی می‌کنند \cite{silvanoIotaTangleCryptocurrency2020}. در نسخه‌های نخست \lr{IOTA}، هر تراکنش جدید دو تراکنش قبلی تأییدنشده (نوک) را با الگوریتم انتخاب نوک مبتنی بر زنجیرهٔ مارکوف مونت‌کارلو\LTRfootnote{\lr{Markov Chain Monte Carlo (MCMC)}} انتخاب و تأیید می‌کند \cite{silvanoIotaTangleCryptocurrency2020}، در حالی که در نسخهٔ دوم پیام‌ها باید به ۲ تا ۸ پیام قبلی ارجاع دهند \cite{IOTATANGLE20, sealeyIOTATangle202022}.

فرایند افزودن تراکنش پنج گام دارد: گره دو تراکنش دیگر را برای تأیید انتخاب می‌کند؛ بررسی می‌کند که این دو در تعارض نباشند و تراکنش‌های متعارض (دوبارخرج‌کردن) را رد می‌کند؛ برای صدور تراکنش معتبر نوعی اثبات کار بسیار سبک‌تر از آن زنجیره بلوکی حل می‌کند؛ تراکنش به شبکه ارسال شده و به یک نوک (تراکنش تأییدنشده) تبدیل می‌شود و منتظر می‌ماند تا وزن انباشته‌اش به آستانهٔ از‌پیش‌تعیین‌شده برسد \cite{silvanoIotaTangleCryptocurrency2020}. از آنجا که کاربران خودشان تراکنش‌ها را تأیید می‌کنند، نیازی به ماینر نیست و هزینهٔ تراکنش تنها شامل هزینهٔ محاسباتی تأیید دو تراکنش دیگر است \cite{silvanoIotaTangleCryptocurrency2020}؛ تراکنش‌ها بدون کارمزد هستند \cite{silvanoIotaTangleCryptocurrency2020, IOTATANGLE20} و این فناوری نسبت به زنجیره بلوکی کارآمدتر از نظر انرژی است \cite{silvanoIotaTangleCryptocurrency2020, sealeyIOTATangle202022, IOTATANGLE20}.

با این حال، نسخه‌های ۱٫۰ (۲۰۱۵) و ۱٫۵ (۲۰۱۷) \lr{IOTA} برای اعتبارسنجی تراکنش و امنیت به یک هماهنگ‌کننده\LTRfootnote{\lr{coordinator}} متمرکز متکی بودند \cite{IOTATANGLE20}؛ این هماهنگ‌کننده هر دو دقیقه یک «نقطهٔ عطف» امضاشده تولید می‌کند و تراکنش‌هایی که مستقیم یا غیرمستقیم به آن ارجاع دارند با اطمینان ۱۰۰ درصدی تأیید می‌شوند \cite{silvanoIotaTangleCryptocurrency2020}. انتقادات مهمی به این سازوکار وارد است: هماهنگ‌کننده به بنیاد \lr{IOTA} اجازه انتخاب اولویت تراکنش‌ها، یخ‌زدن وجوه و نادیده‌گرفتن تراکنش‌ها می‌دهد و نقطهٔ حملهٔ متمرکز است؛ اگر متوقف شود، تأییدها متوقف می‌شوند \cite{silvanoIotaTangleCryptocurrency2020}. بنابراین \lr{IOTA} در نسخه‌های نخست «نیمه‌غیرمتمرکز» بود و بزرگ‌ترین چالش آن حذف هماهنگ‌کننده بود \cite{silvanoIotaTangleCryptocurrency2020}؛ برنامهٔ حذف آن با عنوان «کوردیساید»\LTRfootnote{\lr{coordicide}} در مه ۲۰۱۹ پیشنهاد شد \cite{silvanoIotaTangleCryptocurrency2020}.

نسخهٔ دوم (طراحی‌شده در ۲۰۲۰) هماهنگ‌کنندهٔ متمرکز را حذف و آن را با سازوکار هماهنگی غیرمتمرکزی جایگزین کرد که به گره‌ها اجازه تأیید جمعی تراکنش‌ها و دستیابی به اجماع بدون اتکا به یک موجودیت واحد می‌دهد \cite{IOTATANGLE20}. در سال ۲۰۲۱ نخستین شبکهٔ کاملاً غیرمتمرکز \lr{IOTA} در شبکهٔ توسعه منتشر شد \cite{sealeyIOTATangle202022} و پیاده‌سازی کاری به نام \lr{Nectar}\LTRfootnote{\lr{Nectar}} طبق گزارش‌های بنیاد \lr{IOTA} از تا ۱۰۰۰ تراکنش در ثانیه با زمان تأیید میانگین ۱۰ تا ۱۲ ثانیه پشتیبانی می‌کند؛ این ارقام از منابع بنیاد \lr{IOTA} نقل شده و ارزیابی مستقل نیستند \cite{sealeyIOTATangle202022}. در زمان نگارش منابع، نسخهٔ ۲٫۰ هنوز در حال توسعه بود \cite{IOTATANGLE20}. چرخهٔ اجماع شامل کنترل ازدحام، رأی‌گیری با اجماع احتمالاتی سریع\LTRfootnote{\lr{Fast Probabilistic Consensus (FPC)}} با محافظت مانا، انتخاب نوک از شاخه‌های درست، رشد وزن تأیید و نهایی‌سازی (رسیدن به آستانهٔ معمولاً ۷۰ درصد ظرف ۱۰ ثانیه) و به‌روزرسانیٔ مانا است \cite{IOTATANGLE20}. \lr{FPC} پروتکل رأی‌گیری دودویی بدون رهبر و احتمالاتی است که در آن گره‌ها در هر دور با تعداد محدودی از گره‌های دیگر تعامل می‌کنند \cite{IOTATANGLE20, sealeyIOTATangle202022}؛ آستانه‌های تصادفی توسط کمیتهٔ مولد اعداد تصادفی غیرمتمرکز\LTRfootnote{\lr{decentralized Random Number Generator (dRNG)}} بر اساس مانا تولید می‌شوند \cite{IOTATANGLE20}. شبیه‌سازی \lr{FPC} نشان می‌دهد نرخ توافق و میانگین دور خاتمه تا ۱۰۰۰۰ گره عملاً بدون تغییر می‌ماند و پروتکل تا نقطهٔ غیرواقعی‌ستی که نیمی از شبکه به‌صورت بدخواهانه کنترل شود مقاوم است، هرچند سربار ارتباطی پرس‌وجوهای مستقیم برای برخی پیاده‌سازی‌های \lr{IoT} نگران‌کننده است \cite{sealeyIOTATangle202022}. مانا\LTRfootnote{\lr{mana}} سامانه‌ای شهرت‌مانند برای حفاظت شبکه در برابر حملات سیبل و کسوف\LTRfootnote{\lr{Sybil and Eclipse attacks}} است و دو نوع «مانای دسترسی» و «مانای اجماع» دارد که با گذشت زمان میرا می‌شود \cite{sealeyIOTATangle202022}. در نسخهٔ دوم، اثبات کار تطبیقی برای کنترل نرخ (نه ماینینگ) به‌کار می‌رود: دشواری بر اساس نرخ پیام گره تنظیم می‌شود تا برای گره‌های صادق کم بماند و برای دستگاه‌های محدود \lr{IoT} مفید باشد \cite{sealeyIOTATangle202022}.

\begin{figure}[H]
\centering
\includegraphics[width=0.95\textwidth]{IOTA vs Blockchain.png}
\caption{مقایسه تصویری \lr{IOTA Tangle} با زنجیره بلوکی سنتی}
\label{fig:iota_vs_bc}
\end{figure}

معماری \lr{IOTA Tangle} نسخهٔ دوم از سه لایه تشکیل شده است \cite{IOTATANGLE20, sealeyIOTATangle202022}: لایهٔ شبکه عملیات سطح بایت و همهٔ تبادلات همتا به همتا را مدیریت می‌کند؛ لایهٔ ارتباطات پیام‌های \lr{IOTA} را مدیریت می‌کند و شامل تأییدهای انتخاب نوک، سازوکارهای کنترل نرخ و مدیریت دفتر کل تنگل است؛ و لایهٔ کاربردی محموله‌های پیام را اجرا می‌کند و برای تراکنش ارزشی، اجرای اجماع، انتقال وجه و انتقال مانا (شهرت) را بر عهده دارد \cite{IOTATANGLE20, sealeyIOTATangle202022}. قراردادهای هوشمند از طریق پروتکل قرارداد هوشمند \lr{IOTA}\LTRfootnote{\lr{IOTA Smart Contracts Protocol (ISCP)}} پشتیبانی می‌شوند؛ اجرا به‌صورت خارج‌زنجیره‌ای روی زنجیره‌های فرعی متصل به تنگل اصلی و توسط کمیته‌ای از گره‌ها انجام می‌شود، هزینه‌های پایین و قابل پیش‌بینی دارد، امنیت با افزایش اندازهٔ کمیته بیشتر می‌شود و از ماشین مجازی اتریوم و اجرای مستقیم قراردادهای \lr{Solidity} پشتیبانی می‌کند \cite{IOTATANGLE20, sealeyIOTATangle202022}؛ مالک می‌تواند پاداش متغیری برای کمیته تعیین کند و کارمزد قراردادها به زنجیره، جزئیات قرارداد و مالک بستگی دارد \cite{IOTATANGLE20}.

در حوزهٔ امضا، \lr{IOTA} ۱٫۰ از امضای یک‌بارمصرف وینترنیتز\LTRfootnote{\lr{Winternitz One-Time Signature (W-OTS)}} استفاده می‌کرد که هش‌محور و مقاوم در برابر کامپیوترهای کوانتومی توصیف شده، اما هر بار امضا ۵۰ درصد کلید خصوصی مرتبط با آدرس را افشا می‌کند و استفادهٔ مجدد از آدرس را ناایمن می‌سازد \cite{silvanoIotaTangleCryptocurrency2020}. \lr{IOTA} ۱٫۵ و ۲٫۰ برای بهبود امنیت و استواری به طرح امضای \lr{Ed25519}\LTRfootnote{\lr{Ed25519}} مهاجرت کردند \cite{IOTATANGLE20} که نسخه‌ای از طرح \lr{EdDSA} روی منحنی ادواردز پیچیده با عملکرد خوب و خواص امنیتی استوار است \cite{IOTATANGLE20}. در مقایسهٔ پیاده‌سازی‌شده روی \lr{Raspberry Pi} ۱، اندازهٔ کلید ۳۲ بایت (در مقابل ۲۵۶ بایت)، اندازهٔ امضا ۶۴ بایت (در مقابل ۱۳۰۰ تا ۳۹۰۰ بایت)، زمان تولید کلید ۰٫۱۳۸ ثانیه (در مقابل ۰٫۶۹۶)، زمان امضا ۰٫۱۴۹ (در مقابل ۰٫۳۷۱) و زمان تأیید ۰٫۲۰۷ ثانیه (در مقابل ۰٫۲۵۸) و سطح امنیت ۱۲۸ بیت (در مقابل ۱۱۲ بیت) گزارش شده است \cite{IOTATANGLE20}؛ این اندازه‌گیری‌ها روی سخت‌افزار بسیار ضعیف و با پیاده‌سازی پایتون انجام شده و نمایانگر کارایی در واحد سوارشوندهٔ خودرو نیستند \cite{IOTATANGLE20}. از نظر مقیاس‌پذیری، با افزایش تراکنش‌های \lr{IOTA} شبکه استوارتر شده و تأییدها سریع‌تر انجام می‌شوند \cite{IOTATANGLE20}؛ از آنجا که هر گره هنگام صدور یک پیام باید دو پیام را تأیید کند، با افزایش تعداد تراکنش در ثانیه، زمان تأیید کاهش می‌یابد \cite{sealeyIOTATangle202022}. البته در مرور سیلوانو و مارچلینو، زمان‌های اندازه‌گیری‌شده در مقالات مرورشده از چند ثانیه تا بیش از یک دقیقه گزارش شده و «تأیید سریع» باید با احتیاط بیان شود \cite{silvanoIotaTangleCryptocurrency2020}.

\begin{table}[H]
\centering
\caption{مقایسهٔ \lr{IOTA Tangle} با زنجیره بلوکی}
\label{tab:iota_vs_bc}
\begin{tabular}{p{3.5cm} p{4.5cm} p{4.5cm}}
\toprule
\textbf{ویژگی} & \textbf{زنجیره بلوکی} & \textbf{\lr{IOTA Tangle}} \\
\midrule
ساختار & زنجیره بلوک‌ها (گراف جهت‌دار بدون دور محدود به یک مسیر) & گراف جهت‌دار بدون دور (بدون بلوک) \\
کارمزد & دارد (نوسانی) & بدون کارمزد برای تراکنش‌های پایه \\
تراکنش در ثانیه & محدود (حدود ۷--۱۵) & قابل ارتقا (شبکهٔ \lr{Nectar}: تا ۱۰۰۰) \\
مصرف انرژی & بالا & کمتر \\
ماینر & بلاکچین‌های اثبات کار نیاز دارند & نیاز ندارد (اثبات کار سبک برای هر تراکنش) \\
غیرمتمرکزسازی & بالا & ۱.x: نیمه‌غیرمتمرکز (هماهنگ‌کننده)؛ ۲.۰: غیرمتمرکز \\
قرارداد هوشمند & پشتیبانی & در ۱.x نبود؛ در ۲.۰ از طریق \lr{ISCP} \\
\bottomrule
\end{tabular}
\end{table}

جدول \ref{tab:iota_vs_bc} تفاوت‌های اصلی \lr{IOTA Tangle} و زنجیره بلوکی را خلاصه می‌کند. در شبکهٔ اقتضایی خودرویی نیز کاربردهایی برای \lr{IOTA Tangle} بررسی شده است. طرح \lr{DIVA} یک سامانهٔ شهرت مبتنی بر شناسهٔ غیرمتمرکز\LTRfootnote{\lr{Decentralized Identifier (DID)}} برای انتقال ایمن در شبکه‌های خودرویی با استفاده از \lr{IOTA} است \cite{feraudoDIVADIDbasedReputation2024}: هر خودرو یک \lr{DID} دارد که مرجع مورد اعتماد آن را بر دفتر کل ثبت و برای آن اعتبارنامهٔ تأییدپذیر\LTRfootnote{\lr{Verifiable Credential (VC)}} صادر می‌کند و خودرو با پیوستن ارائهٔ تأییدپذیر\LTRfootnote{\lr{Verifiable Presentation (VP)}} به هر پیام از مالکیت شناسهٔ خود پشتیبانی می‌کند \cite{feraudoDIVADIDbasedReputation2024}. امتیاز شهرت روی گره‌های لبه\LTRfootnote{\lr{edge node}} (نه خودروها) محاسبه می‌شود، با مقایسهٔ پیام اعلان محیطی غیرمتمرکز\LTRfootnote{\lr{Decentralized Environmental Notification Message (DENM)}} با دادهٔ واحد کنارجاده‌ای، پیام‌های آگاهی همیارانه\LTRfootnote{\lr{Cooperative Awareness Message (CAM)}} همان فرستنده و رویدادهای مشابه، و امتیاز نهایی با تراکنش بدون ارزش\LTRfootnote{\lr{zero-value transaction}} در تنگل ذخیره می‌شود \cite{feraudoDIVADIDbasedReputation2024}. دفتر کل مورد استفاده از نوع مجوزدار و مبتنی بر گراف جهت‌دار بدون دور است که فقط مجموعه‌ای محدود از گره‌ها مجاز به انتشار و خواندن تراکنش هستند و تراکنش‌های بدون ارزش نیازی به اعتبارسنجی شرکت‌کنندگان شبکه ندارند \cite{feraudoDIVADIDbasedReputation2024}. در شبیه‌سازی (پنج گره لبه، سناریوی تقاطع اینگولشتات نورد در آلمان با ۲۰ تا ۴۰ درصد منابع مخرب)، تأخیر انتشار شهرت چند میکروثانیه بود و با آستانهٔ میانگین و ضریب \lr{β} بیش از ۰٫۳، دقت حدود ۹۴ درصد با ۳۰ درصد منبع مخرب و ۸۹ درصد با ۴۰ درصد گزارش شده است \cite{feraudoDIVADIDbasedReputation2024}؛ این تأخیر برای ارتباط یک‌به‌یک در شبیه‌سازی اندازه‌گیری شده و هماهنگ‌کنندهٔ \lr{IOTA} را موقت می‌داند \cite{feraudoDIVADIDbasedReputation2024}.

در مرور سیلوانو و مارچلینو، دو کاربرد خودرویی \lr{IOTA} گزارش شده است \cite{silvanoIotaTangleCryptocurrency2020}. طرح \lr{IOTA-VPKI} بر زیرساخت کلید عمومی خودرویی\LTRfootnote{\lr{Vehicular Public Key Infrastructure (VPKI)}} \lr{SECMACE} بنا شده است، بازیگرانی برای ثبت‌نام خودرو و صدور شبه‌نام‌های مورد اعتماد دارد، محرمانگی ثبت و به‌روزرسانی گواهی را با \lr{IOTA} تضمین می‌کند، هر خودرو کلیدی برای دسترسی به کانال پیام‌رسانی احرازشدهٔ پوشیده\LTRfootnote{\lr{Masked Authenticated Messaging (MAM)}} ثبت می‌کند و در برابر حملات ممانعت از سروی مقاوم است \cite{silvanoIotaTangleCryptocurrency2020}. طرح بارتولومئو و همکاران نیز \lr{IOTA} را لایه‌ای امنیتی و حفظ حریم خصوصی برای پروتکل‌های ارتباطی معرفی می‌کند، تأخیر تراکنش کمتری نسبت به زنجیره بلوک‌های عمومی گزارش می‌کند (ارسال تراکنش تا ۶۵۰ میلی‌ثانیه) و تأخیر را برای کاربردهای خودرویی ناچیز می‌داند \cite{silvanoIotaTangleCryptocurrency2020}. همچنین چارچوب هویت \lr{IOTA} روی تنگل بنا شده و در بافت‌هایی مانند حمل‌ونقل هوشمند که موجودیت‌هایی مانند حسگرها و خودروهای هوشمند به هویت دیجیتال نیاز دارند، راه‌حل طبیعی‌ای است و ذخیرهٔ پیام‌ها، از جمله اسناد \lr{DID}، بدون تحمّل کارمزد ممکن است \cite{mazzocca2025survey}. در مقابل، مزایای گراف جهت‌دار بدون دور نسبت به بلاکچین از منظر طرح‌های شهرت خودرویی، مقیاس‌پذیری، تراکنش‌های سریع و کارمزد پایین گزارش شده است \cite{feraudoDIVADIDbasedReputation2024}.
